\documentclass[11pt, letterpaper]{article}

% --- Packages ---
\usepackage[margin=1in]{geometry}
\usepackage{titlesec}
\usepackage{enumitem}
\usepackage{booktabs}
\usepackage{longtable}
\usepackage{hyperref}
\usepackage{xcolor}
\usepackage{fancyhdr}
\usepackage{graphicx}
\usepackage{array}
\usepackage{tabularx}
\usepackage{tcolorbox}
\usepackage[T1]{fontenc}
\usepackage{lmodern}
\usepackage{parskip}

% --- Brand Colors ---
\definecolor{Garnet}{HTML}{73000A}
\definecolor{Black90}{HTML}{363636}
\definecolor{Black70}{HTML}{5C5C5C}
\definecolor{Black50}{HTML}{A2A2A2}
\definecolor{Black30}{HTML}{C7C7C7}
\definecolor{Black10}{HTML}{ECECEC}
\definecolor{Atlantic}{HTML}{466A9F}
\definecolor{Congaree}{HTML}{1F414D}
\definecolor{Rose}{HTML}{CC2E40}
\definecolor{Horseshoe}{HTML}{65780B}
\definecolor{Honeycomb}{HTML}{A49137}

% --- Hyperlink Setup ---
\hypersetup{
    colorlinks=true,
    linkcolor=Garnet,
    urlcolor=Atlantic,
    citecolor=Congaree,
    pdftitle={USC Magellan Scholar Grant -- Comprehensive Applicant Guide},
    pdfauthor={J.C. Vaught},
}

% --- Section Formatting ---
\titleformat{\section}{\Large\bfseries\color{Garnet}}{}{0em}{}[\titlerule]
\titleformat{\subsection}{\large\bfseries\color{Black90}}{}{0em}{}
\titleformat{\subsubsection}{\normalsize\bfseries\color{Black70}}{}{0em}{}

% --- Header/Footer ---
\pagestyle{fancy}
\fancyhf{}
\fancyhead[L]{\small\color{Black70}USC Magellan Scholar Grant Guide}
\fancyhead[R]{\small\color{Black70}2025--2026}
\fancyfoot[C]{\small\color{Black70}\thepage}
\renewcommand{\headrulewidth}{0.4pt}
\renewcommand{\footrulewidth}{0pt}

% --- tcolorbox styles ---
\tcbset{
    sharp corners,
    boxrule=0.5pt,
    left=6pt,
    right=6pt,
    top=4pt,
    bottom=4pt,
}

\newcommand{\tipbox}[2]{%
\begin{tcolorbox}[colback=Black10, colframe=Garnet, title={\textbf{#1}}]
#2
\end{tcolorbox}
}

\newcommand{\warnbox}[1]{%
\begin{tcolorbox}[colback=Black10, colframe=Rose]
\textbf{Important:} #1
\end{tcolorbox}
}

% ==============================================================================
\begin{document}

% --- Title Page ---
\begin{titlepage}
\centering
\vspace*{2cm}
{\Huge\bfseries\color{Garnet} Magellan Scholar Grant\par}
\vspace{0.5cm}
{\LARGE\color{Black90} Comprehensive Applicant \& Mentor Guide\par}
\vspace{1cm}
{\large University of South Carolina\par}
{\large Office of Undergraduate Research\par}
\vspace{1.5cm}
{\large\color{Black70} Academic Year 2025--2026\par}
\vspace{2cm}
{\normalsize Compiled by J.C. Vaught\par}
{\normalsize January 2026\par}
\vfill
{\small\color{Black50} This document is an unofficial compilation of publicly available information from the USC Office of Undergraduate Research. Always consult the official \href{https://sc.edu/about/offices_and_divisions/undergraduate_research/documents/mgs_guidebook.pdf}{Magellan Scholar Guidebook} and \href{https://sc.edu/about/offices_and_divisions/undergraduate_research/funding_opportunities/our_funding/magellan-scholar-award/}{OUR website} for the most current requirements.\par}
\end{titlepage}

\tableofcontents
\newpage

% ==============================================================================
\section{Program Overview}
% ==============================================================================

The \textbf{Magellan Scholar Award} is a competitive research grant administered by the University of South Carolina's \href{https://sc.edu/about/offices_and_divisions/undergraduate_research/}{Office of Undergraduate Research (OUR)}. Established in 2005, the program has funded over 1,000 undergraduate research, scholarly, and creative projects across every discipline---from science, technology, and engineering to theatre, music, and art.

\subsection{Key Facts at a Glance}

\begin{tabularx}{\textwidth}{@{}>{\bfseries}l X@{}}
\toprule
Award Amount & Up to \$2,500 per student \\
Eligibility & Undergraduate, degree-seeking students at \emph{any} USC system campus \\
GPA Requirement & Minimum institutional GPA of 3.200 \\
Disciplines & All disciplines (STEM, social sciences, humanities, arts, business) \\
Grant Type & Research grant (NOT a scholarship; does not affect financial aid) \\
Use of Funds & Salary, materials/supplies, research travel \\
Mentor & Faculty, staff, or community mentor required \\
Award Duration & Approximately 12 months \\
\bottomrule
\end{tabularx}

\warnbox{The Magellan Scholar Grant is only available to \textbf{undergraduate} students. Graduate students are not eligible. For graduate research funding, see the \href{https://sc.edu/about/offices_and_divisions/research_and_grant_development/funding/internal_funding/sparc/}{SPARC Graduate Research Grant} (\$5,000, PhD/MFA/MA Public History only).}

\subsection{Discipline Distribution}

Historically, about \textbf{50\%} of awards go to science and engineering students, \textbf{25\%} to social science students, and \textbf{25\%} to students in the humanities, arts, and business.

\subsection{Campus Eligibility}

The Magellan Scholar Grant is open to students on \textbf{all USC system campuses}:

\begin{itemize}[nosep]
    \item USC Columbia
    \item USC Aiken
    \item USC Beaufort
    \item USC Upstate
    \item USC Lancaster (Palmetto College)
    \item USC Salkehatchie (Palmetto College)
    \item USC Sumter (Palmetto College)
    \item USC Union (Palmetto College)
\end{itemize}

\textbf{Note:} Different campuses may have different submission deadlines. In Fall 2024, Hurricane Helene caused extended deadlines for some campuses. Always verify campus-specific deadlines with OUR.

% ==============================================================================
\section{Eligibility Requirements}
% ==============================================================================

\subsection{Student Eligibility}

\begin{enumerate}[nosep]
    \item Must be a \textbf{USC degree-seeking undergraduate student} (any campus, any discipline, any academic year).
    \item Minimum institutional \textbf{GPA of 3.200}.
    \item Must have an identified \textbf{faculty mentor} willing to collaborate on the project.
    \item Must complete the \textbf{Application Workshop Video} (required once; if completed in a prior semester, does not need to be repeated).
\end{enumerate}

\subsection{Who Cannot Apply}

\begin{itemize}[nosep]
    \item Graduate students (Master's, PhD, professional degree students).
    \item Non-degree-seeking or visiting students.
    \item Students below a 3.200 institutional GPA.
\end{itemize}

\subsection{Repeat Applicants}

Students may \textbf{resubmit proposals} that were not funded in a previous cycle. However:
\begin{itemize}[nosep]
    \item You must follow the \textbf{current semester's guidebook}, forms, and formatting---do not reuse old forms.
    \item Pay special attention to the \textbf{resubmission section} in the guidebook.
    \item Changes occur each application round; it is your responsibility to follow the latest requirements.
\end{itemize}

% ==============================================================================
\section{Deadlines \& Timeline (2025--2026)}
% ==============================================================================

\begin{center}
\begin{tabularx}{\textwidth}{@{}l l l X@{}}
\toprule
\textbf{Cycle} & \textbf{Deadline} & \textbf{Announcement} & \textbf{Grant Period} \\
\midrule
Fall 2025 & Oct 8, 2025 by 5:00 PM & By Dec 15, 2025 & $\sim$Dec 2025 -- Nov 2026 \\
Spring 2026 & Feb 11, 2026 by 5:00 PM & By Apr 17, 2026 & $\sim$Apr 2026 -- Mar 2027 \\
\bottomrule
\end{tabularx}
\end{center}

\warnbox{For the 2025--2026 academic year, a \textbf{large number of applications} is anticipated. This will likely result in a \textbf{lower percentage of funded applications} compared to previous years. Plan accordingly and submit your strongest possible proposal.}

\subsection{Special Deadlines}

\begin{itemize}[nosep]
    \item \textbf{College of Nursing} and \textbf{Arnold School of Public Health} have modified submission procedures and earlier internal deadlines. Contact your respective college coordinator for details.
    \item \textbf{USCeRA approval routing} should be complete at least \textbf{3 business days} before the OUR deadline (allows 2 business days for unit head and dean review).
\end{itemize}

% ==============================================================================
\section{Application Process Step-by-Step}
% ==============================================================================

\subsection{Step 1: Read the Guidebook}

The \href{https://sc.edu/about/offices_and_divisions/undergraduate_research/documents/mgs_guidebook.pdf}{Magellan Scholar Guidebook} is the definitive resource. Read it thoroughly before beginning your proposal. It covers eligibility, formatting, proposal content, forms, and compliance requirements.

\subsection{Step 2: Identify a Faculty Mentor}

Your mentor must be a USC faculty or staff member (or approved community mentor). The mentor will:
\begin{itemize}[nosep]
    \item Collaborate on proposal development.
    \item Submit the application through USCeRA.
    \item Complete the Mentor Collaboration Form.
    \item Authorize all expenditures during the award period.
\end{itemize}

\subsection{Step 3: Watch the Application Workshop Video}

Required for first-time applicants. If you watched it in a previous semester, you do not need to repeat it.

\subsection{Step 4: Develop Your Proposal}

Write a competitive research, scholarly, or creative project proposal in collaboration with your faculty mentor. The key to a successful application is \textbf{how well you describe your project plan}.

\subsection{Step 5: Prepare Budget}

Develop a detailed budget using the current semester's budget form. The maximum request is \$2,500. Funds can cover:
\begin{itemize}[nosep]
    \item Student salary/stipend
    \item Materials and supplies (generally $\leq$10\% of total budget)
    \item Research travel (domestic or international)
\end{itemize}

\subsection{Step 6: Assemble Application Documents}

Combine into \textbf{ONE Word or PDF document}:
\begin{itemize}[nosep]
    \item Project description / proposal
    \item Budget form and justification
    \item Primary Faculty Collaboration Form (mentor completes this)
    \item Secondary Mentor Form(s), if applicable
    \item Non-USC Co-Mentor Letter, if applicable
\end{itemize}

\textbf{Group Proposals:} Name the file \texttt{StudentLastName1\_StudentLastName2}. The USCeRA title should follow the format: \texttt{Magellan-StudentLastName1\_StudentLastName2-ProjectTitle}.

\subsection{Step 7: Submit via USCeRA}

The mentor submits the full application through the \href{https://sam.research.sc.edu/uscera/}{USCeRA system}. Steps:
\begin{enumerate}[nosep]
    \item Log into USCeRA.
    \item Complete the Proposal/Award Processing (PAP) form.
    \item Upload all proposal components.
    \item Complete supplemental web-based forms as prompted.
    \item Route the proposal via the ``Start Approval Process'' button.
    \item Proposal routes: Chair $\rightarrow$ Dean $\rightarrow$ SAM for approval.
\end{enumerate}

\tipbox{USCeRA Resources}{
\begin{itemize}[nosep]
    \item \href{https://sc.edu/about/offices_and_divisions/undergraduate_research/documents/mgs_uscera-guide.pdf}{Magellan Scholar USCeRA Submission Guide}
    \item \href{https://sam.research.sc.edu/pdf/USCERA_Userguide.pdf}{USCeRA User Guide}
    \item Technical support: Debbie Kassianos, \href{mailto:kassiano@sc.edu}{kassiano@sc.edu}, 803-777-6421
\end{itemize}
}

% ==============================================================================
\section{Proposal Writing Guide}
% ==============================================================================

\subsection{Proposal Structure}

While specific sections vary by discipline, a strong Magellan Scholar proposal typically includes:

\begin{enumerate}[nosep]
    \item \textbf{Introduction / Background} --- Context, literature review, and rationale for the project.
    \item \textbf{Research Question / Creative Objective} --- Clear statement of what you aim to investigate or create.
    \item \textbf{Significance} --- Why does this project matter? What contribution does it make?
    \item \textbf{Methodology / Project Plan} --- Detailed description of your approach, methods, techniques.
    \item \textbf{Timeline} --- Realistic schedule for completing the project within the grant period.
    \item \textbf{Expected Outcomes} --- What you expect to find, produce, or learn.
    \item \textbf{Budget Justification} --- Explanation of how requested funds will be used.
    \item \textbf{References / Bibliography} --- Cited works supporting your proposal.
\end{enumerate}

\subsection{Review Rubric}

Proposals are evaluated by reviewers using a formal rubric. The publicly available rubric (Fall 2019 revision) assesses:

\begin{center}
\begin{tabularx}{\textwidth}{@{}l X l@{}}
\toprule
\textbf{Criterion} & \textbf{Description} & \textbf{Max Points} \\
\midrule
Writing Style & Clear, persuasive, logical; well organized with few errors & 32 \\
Project Description & Quality of research question, methodology, timeline & Varies \\
Significance & Contribution to the field or discipline & Varies \\
Feasibility & Can the project realistically be completed? & Varies \\
Mentor Collaboration & Quality of mentor--student plan & Varies \\
\bottomrule
\end{tabularx}
\end{center}

\textbf{Full rubric:} \href{https://www.sc.edu/about/offices_and_divisions/undergraduate_research/documents/mgs_review_rubric.pdf}{mgs\_review\_rubric.pdf}

\tipbox{Tips for a Strong Proposal}{
\begin{enumerate}[nosep]
    \item \textbf{Start early.} Set up a timeline so your mentor can review multiple drafts.
    \item \textbf{Be specific.} Vague methodology is the most common weakness.
    \item \textbf{Write clearly.} Reviewers come from many disciplines; avoid excessive jargon.
    \item \textbf{Justify your budget.} Every line item should connect to a project activity.
    \item \textbf{Show significance.} Explain why your project matters beyond your own interest.
    \item \textbf{Use the rubric.} Score your own proposal against the review criteria before submitting.
    \item \textbf{Proofread.} Grammatical errors and typos undermine credibility.
\end{enumerate}
}

% ==============================================================================
\section{Budget Guidelines}
% ==============================================================================

\subsection{Maximum Award}

\$2,500 per student. You do not have to request the full amount. Group proposals receive up to \$2,500 per student (e.g., a group of 3 could receive up to \$7,500 total).

\subsection{Allowable Expenses}

\begin{itemize}[nosep]
    \item \textbf{Student salary / stipend}
    \item \textbf{Materials and supplies} (generally $\leq$10\% of total)
    \item \textbf{Research travel} (domestic or international; must comply with USC travel policies)
\end{itemize}

\subsection{Prohibited Expenses}

\begin{itemize}[nosep]
    \item Mentor compensation (salary, benefits, or travel)
    \item Graduate student compensation (salary, benefits, or travel)
    \item Other faculty expenses
    \item Tuition for study abroad programs (room/board \emph{may} be covered)
    \item Conference travel (use the \href{https://sc.edu/about/offices_and_divisions/undergraduate_research/funding_opportunities/our_funding/magellan-research-application/}{Magellan Voyager Award} instead)
\end{itemize}

\subsection{How Funds Are Managed}

\begin{itemize}[nosep]
    \item Awards are transferred by the Controller's Office into internal fund ``N'' accounts in the mentor's department.
    \item Setup occurs approximately 3--4 weeks after award announcement.
    \item The \textbf{mentor authorizes all expenditures}.
    \item The department business manager processes transactions.
    \item Funds may only be used for the specific awarded student(s) and project.
\end{itemize}

% ==============================================================================
\section{Group Projects}
% ==============================================================================

Group proposals are permitted under these rules:

\begin{itemize}[nosep]
    \item Maximum of \textbf{3 students per mentor} (one group proposal, or one group plus one individual).
    \item A mentor may not be PI/co-PI on more than \textbf{2 applications per round}.
    \item Mentors with a currently funded group project through Magellan \textbf{may not submit another group proposal}.
    \item Group participation must be \textbf{justified} in both the methods section and the mentor collaboration form.
    \item Each student in a group receives up to \$2,500.
\end{itemize}

% ==============================================================================
\section{Post-Award Requirements}
% ==============================================================================

Once awarded, Magellan Scholars must fulfill the following:

\subsection{CITI Training}

Complete the \textbf{CITI Responsible Conduct of Research (RCR)} online training. Access through the \href{https://about.citiprogram.org/}{CITI Program website}. Two module tracks:
\begin{itemize}[nosep]
    \item \textbf{Social \& Behavioral Research} --- for social/behavioral projects.
    \item \textbf{Biomedical} --- for biomedical projects.
\end{itemize}

\tipbox{CITI Training Resources}{
\begin{itemize}[nosep]
    \item \href{https://sc.edu/about/offices_and_divisions/undergraduate_research/documents/cititraininghowto.pdf}{CITI Training How-To Guide (USC OUR)}
    \item \href{https://sc.edu/about/offices_and_divisions/research_compliance/documents/citi_faq.pdf}{CITI Registration FAQ}
    \item Contact: Office of Research Compliance, 803-576-7276
\end{itemize}
}

\subsection{Project Registration}

Register your project at the \textbf{beginning of each semester} of involvement, and re-register every semester throughout the grant period.

\subsection{Compliance}

\subsubsection{Human Subjects (IRB)}
\begin{itemize}[nosep]
    \item Undergraduate Magellan research is typically \emph{not} required to undergo formal IRB review.
    \item However, students involved in a mentor's IRB-approved project must be added to the existing protocol.
    \item CITI Human Subjects training is recommended for all projects involving people (interviews, surveys, personal information).
    \item Contact: IRB Coordinator, \href{mailto:Newtonla@mailbox.sc.edu}{Newtonla@mailbox.sc.edu}, 803-576-7276.
\end{itemize}

\subsubsection{Animal Subjects (IACUC)}
\begin{itemize}[nosep]
    \item IACUC approval is \textbf{not} required before proposal submission.
    \item Students \textbf{may not begin} animal research until IACUC approval is obtained.
    \item The proposal must include a statement that the mentor is submitting for Animal Use approval.
    \item Contact: Dept.\ of Animal Resources, \href{mailto:iacuc@mailbox.sc.edu}{iacuc@mailbox.sc.edu}, 803-777-8106.
\end{itemize}

\subsubsection{International Research / Travel Abroad}
\begin{itemize}[nosep]
    \item All international travel must comply with USC policies \textbf{and} the \href{https://sc.edu/about/offices_and_divisions/education_abroad/}{Education Abroad Office}.
    \item Students must register with the Education Abroad Office.
    \item USC-sponsored travel health insurance is \textbf{required}.
    \item All travel must be completed \textbf{before graduation}.
    \item Magellan funds will \textbf{not} pay for study abroad tuition (room/board may be covered).
    \item Contact: \href{mailto:edabroad@sc.edu}{edabroad@sc.edu}, 803-777-7557.
\end{itemize}

\subsection{No Extensions}

Extensions past the stated end date or after the student graduates are \textbf{not permitted}.

% ==============================================================================
\section{For Faculty Mentors / PIs}
% ==============================================================================

\subsection{Mentor Responsibilities}

\begin{enumerate}[nosep]
    \item Collaborate with the student on proposal development.
    \item Complete the \textbf{Mentor Collaboration Form} detailing your role and expertise.
    \item Submit the full application via \textbf{USCeRA}.
    \item \textbf{Authorize all expenditures} throughout the award.
    \item Ensure ethical standards and regulatory compliance.
    \item Review and approve final reports before submission.
    \item Oversee student training and skill development.
\end{enumerate}

\subsection{Mentor Limits}

\begin{itemize}[nosep]
    \item No more than \textbf{2 applications per round} as PI/co-PI.
    \item Maximum of \textbf{3 students} (one group or one group + one individual).
    \item Cannot submit a new group proposal if a group project is currently funded.
\end{itemize}

\subsection{Sabbatical / Extended Leave}

Mentors anticipating sabbaticals or extended leave during the project period must:
\begin{itemize}[nosep]
    \item Identify a \textbf{secondary faculty mentor} to supervise during the absence.
    \item Alternatively, consider delaying the proposal until after return from leave.
\end{itemize}

\subsection{Non-USC Co-Mentors}

Non-USC co-mentors may provide a \textbf{letter of support} in lieu of the standard mentor collaboration form.

% ==============================================================================
\section{Examples of Funded Projects}
% ==============================================================================

The following are publicly documented examples of Magellan Scholar projects across disciplines, illustrating the breadth of the program.

\subsection{Biological Sciences}

\begin{longtable}{@{}p{3.5cm} p{7cm} p{3.5cm}@{}}
\toprule
\textbf{Student} & \textbf{Project Title} & \textbf{Mentor} \\
\midrule
\endhead
Anna Bazell & Functional Analysis of STN1 Winged Helix Domains & Dr.\ Jason Stewart \\
William Caspino & Germination Success Following Chronic Radiation Exposure & Dr.\ Tim Mousseau \\
Dylan Davis & Plasticity of Mixotrophy: Molecular and Ecological Responses in \emph{Cryptomonas ovata} & Dr.\ Jeff Dudycha \\
Kaitlyn Dirr & Seasonal Use of Gopher Tortoise Burrows in Coastal Beach Dunes & Dr.\ Josh Stone \\
Kirsten Fisher & Perception of Green Leaf Volatiles by Plant Cells & Dr.\ Johannes Stratmann \\
Megan Glover & Effects of High-Fat Diets on Adipose Tissue and Immune Function & Dr.\ Alissa Armstrong \\
Benjamin Hurley & DNA Barcoding to Assess Dietary Composition of Highly Migratory Species in the Northwest Atlantic & Dr.\ Joe Quattro \\
Catherine Alexander & Quantifying Urban Artificial Light on Pollination Success of \emph{Glandularia bipinnatifida} & (Faculty mentor) \\
\bottomrule
\end{longtable}

\subsection{Psychology \& Behavioral Sciences}

\begin{longtable}{@{}p{3.5cm} p{9cm} p{2cm}@{}}
\toprule
\textbf{Student} & \textbf{Project Title} & \textbf{Year} \\
\midrule
\endhead
Ella Ahrens & Premature Language Decline in Women with the FMR1 Premutation & --- \\
Sydney Burrell & Mother-Child Synchrony as a Predictor of Problem Behavior and Autism-Related Deficits in Individuals with Fragile X Syndrome & --- \\
\bottomrule
\end{longtable}

\subsection{Education}

\begin{longtable}{@{}p{5cm} p{7.5cm} p{2cm}@{}}
\toprule
\textbf{Students} & \textbf{Project Title} & \textbf{Campus} \\
\midrule
\endhead
Adam Taylor Eubanks, Deanna McCord, Skyla Watson & Examining the Perceived Self-Efficacy of Inclusive Special Area Teachers in K-12 Classrooms & USC Aiken \\
\bottomrule
\end{longtable}

\textbf{Note:} This was a group proposal; the three students received a combined \$7,500.

\subsection{History \& Humanities}

\begin{longtable}{@{}p{3.5cm} p{7.5cm} p{3.5cm}@{}}
\toprule
\textbf{Student} & \textbf{Project Title} & \textbf{Details} \\
\midrule
\endhead
Monica Bowman & Richard Howell Gleaves: Reconstruction-Era Politician in South Carolina & 1,000th Magellan Scholar; mentored by Dr.\ Melissa Cooper \\
\bottomrule
\end{longtable}

\subsection{International Business}

\begin{longtable}{@{}p{3.5cm} p{7.5cm} p{3.5cm}@{}}
\toprule
\textbf{Student} & \textbf{Project Details} & \textbf{Mentor} \\
\midrule
\endhead
Joy Gannaway & Immersive study in Mauritius (international business research) & Prof.\ Lite Nartey \\
\bottomrule
\end{longtable}

\subsection{Medical / Public Health}

Dr.\ Anna Hoppmann, now a pediatric oncologist, received a Magellan Scholar award as an anthropology major to conduct community-based participatory research on the need for medical interpretation in clinical settings---an example of how Magellan research shapes long-term careers.

\subsection{Broader Discipline Coverage}

Past Magellan Scholar projects have spanned: opera direction, environmental health, marine biology, Czech history, community health outreach, chemical engineering, electrical engineering, computing, and biomedical engineering, among many others.

% ==============================================================================
\section{Related USC Funding Opportunities}
% ==============================================================================

\subsection{Undergraduate Research Mini-Grant}

\begin{tabularx}{\textwidth}{@{}>{\bfseries}l X@{}}
\toprule
Amount & Up to \$1,000 per student \\
Disciplines & All \\
Use & Salary, materials, research travel (NOT conference travel) \\
Fall 2025 Deadline & October 21, 2025, 5:00 PM \\
Spring 2026 Deadline & March 24, 2026, 5:00 PM \\
GPA Requirement & Not specified (lower barrier than Magellan Scholar) \\
\bottomrule
\end{tabularx}

\subsection{Magellan Voyager Award}

\begin{tabularx}{\textwidth}{@{}>{\bfseries}l X@{}}
\toprule
Amount & Up to \$500 \\
Purpose & Travel to conferences to present research \\
Note & Research does NOT need to be OUR-sponsored \\
\bottomrule
\end{tabularx}

\subsection{Conference Presentation Grant}

Up to \$500 for travel to conferences, symposiums, and meetings to present posters, give talks, or perform.

\subsection{SPARC Graduate Research Grant (Graduate Only)}

\begin{tabularx}{\textwidth}{@{}>{\bfseries}l X@{}}
\toprule
Amount & Up to \$5,000 \\
Eligibility & PhD (2nd year+), MFA, MA Public History only \\
Website & \href{https://sc.edu/about/offices_and_divisions/research_and_grant_development/funding/internal_funding/sparc/}{go.sc.edu/SPARC} \\
\bottomrule
\end{tabularx}

\subsection{NSF REU Programs at USC}

\begin{longtable}{@{}p{5cm} p{5cm} p{4cm}@{}}
\toprule
\textbf{Program} & \textbf{Focus} & \textbf{Details} \\
\midrule
\endhead
Physics \& Astronomy REU & Nuclear/particle physics, nanotech, modeling & 10 weeks, 8 participants \\
Criminal Justice REU (DCJS) & Race, class, gender in CJ system & 10 weeks, 9 participants \\
Chemical Engineering REU & Nanostructured materials for energy & Summer \\
Quantum Information Science REU & Quantum computing & Summer \\
Mathematics REU & Mathematical Foundations of Data Analysis & NSF RTG funded \\
\bottomrule
\end{longtable}

\subsection{Funding Strategy Summary}

\begin{center}
\begin{tabularx}{\textwidth}{@{}l l X@{}}
\toprule
\textbf{Tier} & \textbf{Amount} & \textbf{Programs} \\
\midrule
1 --- Primary & \$2,500 & Magellan Scholar Grant \\
2 --- Supplemental & \$1,000 & Undergraduate Research Mini-Grant \\
3 --- Travel & \$500 & Magellan Voyager Award; Conference Presentation Grant \\
4 --- Summer & Varies & NSF REU Programs (typically \$5,000+ with stipend) \\
\bottomrule
\end{tabularx}
\end{center}

These can be combined---e.g., a Magellan Scholar could also apply for a Voyager Award to present findings at a conference.

% ==============================================================================
\section{Research Presentation Opportunities}
% ==============================================================================

\subsection{Discover USC}

System-wide showcase of student research, scholarship, and leadership. Multiple presentation formats available. Open to all USC campuses.

\subsection{Discovery Day Abstracts}

Historical abstracts (2004--2014) are archived on \href{https://scholarcommons.sc.edu/discoveryday/}{Scholar Commons}. Browse by discipline or author to see past project examples.

\subsection{Summer Research Symposium}

Poster presentation event for summer research projects across all disciplines.

\subsection{Carolina CrossTalk}

Student-run publication documenting ``the personal story behind research projects.''

\subsection{Regional Campus Events}

\begin{itemize}[nosep]
    \item \textbf{USC Aiken} --- Scholar Showcase (mid-April)
    \item \textbf{USC Beaufort} --- Student Research and Scholarship Day (mid-to-late April)
    \item \textbf{USC Upstate} --- Research Symposium (late March to early April)
\end{itemize}

\warnbox{Before sharing research publicly, verify with your mentor whether your work has commercialization potential. It may fall under intellectual property protections and require clearance.}

% ==============================================================================
\section{Key Contacts}
% ==============================================================================

\begin{center}
\begin{tabularx}{\textwidth}{@{}l X l@{}}
\toprule
\textbf{Office / Contact} & \textbf{Role} & \textbf{Email / Phone} \\
\midrule
Office of Undergraduate Research & General support, applications & \href{mailto:our@sc.edu}{our@sc.edu}, 803-777-1141 \\
 & & Legare College, Suite 120 \\
Debbie Kassianos (SAM) & USCeRA technical support & \href{mailto:kassiano@sc.edu}{kassiano@sc.edu}, 803-777-6421 \\
Office of Research Compliance & CITI training, IRB & 803-576-7276 \\
IRB Coordinator & Human subjects & \href{mailto:Newtonla@mailbox.sc.edu}{Newtonla@mailbox.sc.edu} \\
Dept.\ of Animal Resources & IACUC / animal subjects & \href{mailto:iacuc@mailbox.sc.edu}{iacuc@mailbox.sc.edu}, 803-777-8106 \\
Education Abroad Office & International travel compliance & \href{mailto:edabroad@sc.edu}{edabroad@sc.edu}, 803-777-7557 \\
USC Financial \& Business Svcs & N-account fund transfers & 803-821-1900 \\
\bottomrule
\end{tabularx}
\end{center}

% ==============================================================================
\section{Frequently Asked Questions}
% ==============================================================================

\begin{description}[style=nextline]

\item[Is the Magellan Scholar Grant a scholarship?]
No. It is a research grant. It does \textbf{not} impact your financial aid.

\item[Can graduate students apply?]
No. The Magellan Scholar Grant is exclusively for undergraduate students.

\item[Can I apply from any USC campus?]
Yes. All USC system campuses are eligible. Deadlines may vary by campus.

\item[Can I reapply if my proposal was not funded?]
Yes. Follow the current semester's guidebook and forms. Do not reuse old forms.

\item[Do I have to request the full \$2,500?]
No. Request only what your project needs.

\item[Can I do a group project?]
Yes. Up to 3 students per mentor. Group participation must be justified in the proposal. Each student can receive up to \$2,500.

\item[What if my mentor goes on sabbatical?]
You must identify a secondary mentor to supervise during the absence, or delay your proposal.

\item[Can I use funds for conference travel?]
No. Conference travel is funded through the Magellan Voyager Award (\$500). Magellan Scholar funds are for research-related expenses only.

\item[Do I need IRB approval before applying?]
No. IRB/IACUC approval is not required before proposal submission. However, you must obtain approval before beginning relevant research activities.

\item[Can I do research abroad?]
Yes. You must comply with the Education Abroad Office requirements, register your travel, and hold USC-sponsored travel health insurance. All travel must be completed before graduation.

\end{description}

% ==============================================================================
\section{Sources \& Further Reading}
% ==============================================================================

\begin{enumerate}[nosep]
    \item \href{https://sc.edu/about/offices_and_divisions/undergraduate_research/funding_opportunities/our_funding/magellan-scholar-award/}{Magellan Scholar Grant --- USC OUR}
    \item \href{https://sc.edu/about/offices_and_divisions/undergraduate_research/documents/mgs_guidebook.pdf}{Magellan Scholar Guidebook (PDF)}
    \item \href{https://www.sc.edu/about/offices_and_divisions/undergraduate_research/documents/mgs_review_rubric.pdf}{Proposal Review Rubric (PDF)}
    \item \href{https://sc.edu/about/offices_and_divisions/undergraduate_research/apply_for_funding/our_funding/magellan-scholar-award/applying/index.php}{Applying for the Magellan Scholar Grant}
    \item \href{https://sc.edu/about/offices_and_divisions/undergraduate_research/apply_for_funding/our_funding/magellan-scholar-award/award-requirements/index.php}{Award Requirements}
    \item \href{https://sc.edu/about/offices_and_divisions/undergraduate_research/apply_for_funding/our_funding/magellan-scholar-award/using-your-money/index.php}{Using Your Money}
    \item \href{https://sc.edu/about/offices_and_divisions/undergraduate_research/apply_for_funding/our_funding/magellan-scholar-award/for-mentors/index.php}{For Mentors}
    \item \href{https://sc.edu/about/offices_and_divisions/undergraduate_research/documents/mgs_uscera-guide.pdf}{USCeRA Submission Guide}
    \item \href{https://sc.edu/about/offices_and_divisions/undergraduate_research/documents/cititraininghowto.pdf}{CITI Training How-To Guide}
    \item \href{https://sc.edu/about/offices_and_divisions/research_compliance/irb/student_research.php}{Student Research --- Office of Research Compliance}
    \item \href{https://scholarcommons.sc.edu/discoveryday/}{Discovery Day Abstracts --- Scholar Commons}
    \item \href{https://sc.edu/uofsc/posts/2014/02_julie_morris_undergraduate_research_1000_magellan.php}{Helping Guide a Thousand Magellan Voyages (USC News)}
    \item \href{https://www.sc.edu/about/offices_and_divisions/research/news_and_pubs/news/2017/20170921_magellan_programs.php}{Magellan Programs Support Undergraduate Research (USC VPR)}
\end{enumerate}

\end{document}
